\section*{Capítulo \ref{ch:g6:density-floating} --- Densidade: por que as coisas flutuam}

\begin{solution}{exo:g6:density-floating:1}
A \omterm{def:g3:measuring-mass:mass}{massa} de um \omterm{def:g6:mass-and-volume:volume}{centímetro cúbico} da substância, em \omterm{def:g3:measuring-mass:units}{gramas}; a da água é
\qty{1}{g} por \unit{cm^3}. A regra: menos denso que a água \omterm{def:g1:floating-and-sinking:words}{flutua},
mais denso \omterm{def:g1:floating-and-sinking:words}{afunda}.
\end{solution}

\begin{solution}{exo:g6:density-floating:2}
O carvalho \omterm{def:g1:floating-and-sinking:words}{flutua}; a pedra \omterm{def:g1:floating-and-sinking:words}{afunda}; a cortiça \omterm{def:g1:floating-and-sinking:words}{flutua}; o ferro \omterm{def:g1:floating-and-sinking:words}{afunda};
o gelo \omterm{def:g1:floating-and-sinking:words}{flutua} --- por um fio.
\end{solution}

\begin{solution}{exo:g6:density-floating:3}
$120 \div 80 = 1.5$: \omterm{def:g6:density-floating:density}{densidade} $1.5$ --- mais denso que a água, ele
\omterm{def:g1:floating-and-sinking:words}{afunda}.
\end{solution}

\begin{solution}{exo:g6:density-floating:4}
$158 \div 20 = 7.9$: a assinatura do ferro.
\end{solution}

\begin{solution}{exo:g6:density-floating:5}
A \omterm{def:g6:density-floating:density}{densidade} pertence à substância, e não ao pedaço: cada \omterm{def:g6:mass-and-volume:volume}{centímetro
cúbico} de carvalho carrega a mesma cota de \omterm{def:g3:measuring-mass:units}{gramas}, seja lasca ou
viga. É isso que faz da \omterm{def:g6:density-floating:density}{densidade} uma carteira de identidade para
substâncias desconhecidas.
\end{solution}

\begin{solution}{exo:g6:density-floating:6}
Os \omterm{def:g3:solids-liquids-gases:liquid}{líquidos} se empilham por \omterm{def:g6:density-floating:density}{densidade}, com o mais denso no fundo: mel
$1.4$ abaixo da água $1.0$, abaixo do óleo $0.9$. A uva ($1.05$)
\omterm{def:g1:floating-and-sinking:words}{afunda} pelo óleo e pela água e descansa no piso de mel --- mais densa
que os dois andares de cima, menos densa que o mel.
\end{solution}

\begin{solution}{exo:g6:density-floating:7}
A maioria dos \omterm{def:g3:solids-liquids-gases:solid}{sólidos} é mais densa que o próprio \omterm{def:g3:solids-liquids-gases:liquid}{líquido} e \omterm{def:g1:floating-and-sinking:words}{afunda}
nele. O gelo \omterm{def:g1:floating-and-sinking:words}{flutua} porque a água incha ao \omterm{def:g2:ice-water-steam:meltfreeze}{congelar} --- o mesmo
material num décimo a mais de espaço dá \omterm{def:g6:density-floating:density}{densidade} $0.9 < 1$.
\end{solution}

\begin{solution}{exo:g6:density-floating:8}
A madeira assina cerca de $0.6$: cada \omterm{def:g6:mass-and-volume:volume}{centímetro cúbico} do tronco
oferece menos que \qty{1}{g} de água, então o tronco \omterm{def:g1:floating-and-sinking:words}{flutua} por maior
que seja. O bronze assina perto de $9$: cada \omterm{def:g6:mass-and-volume:volume}{centímetro cúbico} da
moeda oferece mais que a água, então ela \omterm{def:g1:floating-and-sinking:words}{afunda} por menor que seja. O
tamanho nunca importou; a \omterm{def:g3:measuring-mass:mass}{massa} por espaço sempre importou.
\end{solution}

\begin{solution}{exo:g6:density-floating:9}
$\qty{1}{L} = \qty{1000}{cm^3}$, então $800 \div 1000 = 0.8$ \omterm{def:g3:measuring-mass:units}{grama}
por \omterm{def:g6:mass-and-volume:volume}{centímetro cúbico}. O gelo, com $0.9$, é \emph{mais denso} que
esse \omterm{def:g3:solids-liquids-gases:liquid}{líquido}: o cubo de gelo \omterm{def:g1:floating-and-sinking:words}{afunda} nele.
\end{solution}

\begin{solution}{exo:g6:density-floating:10}
O pacote que \omterm{def:g1:floating-and-sinking:words}{flutua} é o casco inteiro: a casca de aço \emph{mais} o
enorme \omterm{def:g6:mass-and-volume:volume}{volume} de \omterm{def:g2:air-around-us:air}{ar} que ela encerra. Na média de todo esse espaço, a
\omterm{def:g6:density-floating:density}{densidade} do pacote cai abaixo de $1$, e a regra manda \omterm{def:g1:floating-and-sinking:words}{flutuar}. Um
rombo deixa a água substituir o \omterm{def:g2:air-around-us:air}{ar}: a \omterm{def:g6:density-floating:density}{densidade} do pacote sobe na
direção da do próprio aço, passa de $1$ --- e a regra, impassível,
manda \omterm{def:g1:floating-and-sinking:words}{afundar}.
\end{solution}

\begin{solution}{exo:g6:density-floating:11}
O juiz mudou: contra o $1.03$ da água do mar, o corpo de um nadador,
perto de $1$, fica confortavelmente do lado de quem \omterm{def:g1:floating-and-sinking:words}{flutua}, e por
isso o mar carrega você mais alto. Entrando na água doce ($1.0$), o
navio perde essa margem e \omterm{def:g1:floating-and-sinking:words}{afunda} mais: a linha d'água sobe pelo
casco.
\end{solution}

\begin{solution}{exo:g6:density-floating:12}
Com metade submersa, a esfera fica exatamente tão fundo na água
quanto deve ficar um pacote com metade da \omterm{def:g6:density-floating:density}{densidade} da água: a casca
mais o \omterm{def:g2:air-around-us:air}{ar} dão na média cerca de $0.5$ \omterm{def:g3:measuring-mass:units}{grama} por \omterm{def:g6:mass-and-volume:volume}{centímetro cúbico} ---
o \omterm{def:g2:air-around-us:air}{ar} escondido fazendo quase todo o trabalho de aliviar. Conforme a
água vaza para dentro, ela substitui o \omterm{def:g2:air-around-us:air}{ar}, a \omterm{def:g6:density-floating:density}{densidade} média do
pacote sobe e a esfera \omterm{def:g1:floating-and-sinking:words}{afunda} cada vez mais; no instante em que a
média passa do $1.0$ da água, a regra da flutuação troca de lado e a
esfera vai ao fundo.
\end{solution}

\begin{solution}{pb:g6:density-floating:1}
\textbf{1.} Nada: trocar ouro por uma \emph{\omterm{def:g3:measuring-mass:mass}{massa}} igual de prata
deixa a balança feliz --- a \omterm{def:g3:measuring-mass:mass}{massa} sozinha não enxerga a
substituição.
\textbf{2.} $1930 \div 19.3 = \qty{100}{cm^3}$.
\textbf{3.} \qty{130}{cm^3} --- o transbordamento é o \omterm{def:g6:mass-and-volume:volume}{volume} da
coroa, por deslocamento.
\textbf{4.} A coroa ocupa \qty{30}{cm^3} a mais de espaço do que
ouro puro da \omterm{def:g3:measuring-mass:mass}{massa} dela deveria ocupar: alguma coisa mais volumosa
que o ouro se esconde lá dentro.
\textbf{5.} $1930 \div 130 \approx 14.8$.
\textbf{6.} Ouro $19.3$ --- coroa $14.8$ --- prata $10.5$: a coroa
cai entre as duas assinaturas.
\textbf{7.} ``Majestade, cada \omterm{def:g6:mass-and-volume:volume}{centímetro cúbico} de ouro puro carrega
\qty{19.3}{g}; a sua coroa carrega apenas uns $15$ --- de ouro puro
ela não é.''
\textbf{8.} A faixa dá \omterm{def:g6:density-floating:density}{densidades} de $1900 \div 130 \approx 14.6$ até
$1960 \div 130 \approx 15.1$ --- nem perto de $19.3$: nenhum erro
plausível de balança salva o ourives.
\textbf{9.} A prata é menos densa: cada \omterm{def:g3:measuring-mass:units}{grama} de ouro roubado foi
devolvido como um \omterm{def:g3:measuring-mass:units}{grama} de prata que precisa de \emph{mais espaço}
--- a troca \omterm{def:g3:measuring-mass:units}{grama} por \omterm{def:g3:measuring-mass:units}{grama} preserva a \omterm{def:g3:measuring-mass:mass}{massa} e infla o \omterm{def:g6:mass-and-volume:volume}{volume}.
\textbf{10.} Parte de ouro: $965 \div 19.3 = \qty{50.0}{cm^3}$;
parte de prata: $965 \div 10.5 \approx \qty{91.9}{cm^3}$.
\textbf{11.} Total $\approx \qty{142}{cm^3}$ --- mais que os $130$
medidos: o palpite meio a meio usou prata demais; o roubo verdadeiro
foi menor (o \omterm{def:g6:mass-and-volume:volume}{volume} medido fica entre $100$ e $142$).
\textbf{12.} ``Provado: a coroa não é de ouro puro. Método: a \omterm{def:g3:measuring-mass:mass}{massa}
dela repartida pelo \omterm{def:g6:mass-and-volume:volume}{volume} medido com água dá uma \omterm{def:g6:density-floating:density}{densidade} muito
abaixo da assinatura imutável do ouro. E a coroa nunca precisou nem
de um arranhão: a água leu o segredo dela por fora --- esse é o
triunfo.''
\end{solution}
